% problem-000215 8 亚油酸酯自由基氧化
\section*{8. 亚油酸酯自由基氧化}
脂类物质的氧化是油脂保存过程中⼴泛存在的问题。亚油酸⼄酯（记为RH）的氧化过程属于⾃由基氧化反应，其中H代表最可能发⽣氧化反应的氢原⼦。反应⽅程式如下：

$
\mathrm{RH} + \mathrm{O} _ {2} \rightarrow \mathrm{ROOH}
$

\subsection*{8-1}
亚油酸的学名是(9�,12�)-9,12-⼗⼋碳⼆烯酸，画出亚油酸⼄酯的结构式，并⽤“\*”标记最可能发⽣⾃由基氧化反应的氢原⼦。

\subsection*{8-2}
实验测得反应的速率⽅程如下：

$
r = - \frac {\mathrm{d} p (\mathrm{O} _ {2})}{\mathrm{d} t} = \frac {a [ \mathrm{ROOH} ] [ \mathrm{RH} ]}{1 + b \frac {[ \mathrm{RH} ]}{p [ \mathrm{O} _ {2} ]}}\tag{1}
$

其中， $a _ { \mathrm { \scriptscriptstyle 3 } }$ 、 $b _ { \scriptscriptstyle * }$ 是与浓度⽆关的常数。该反应的基元反应步骤如下表所⽰，其中， $k _ { 1 }$ 到 $k _ { 5 }$ 分别为相应基元反应的速率常数，链传递步骤是整个反应的决速步。氧⽓分 $\vec { \underline { { \underline { { \tau } } } } } p ( \mathrm { O } _ { 2 } )$ 较⾼时主要按照①、②、④步骤进⾏； $p ( \mathrm { O } _ { 2 } )$ 较低时主要按照①、③、⑤步骤进⾏。据此，分别在两种情况下推导反应的速率⽅程，并解释在 $\dot { \bf \nabla } p ( { \bf O } _ { 2 } )$ 较⾼和较低的情况下(1)式如何近似得出这两种速率⽅程。

\textit{（原卷此处含表格/仪器试剂表，详见 PDF 或 MD 源文件。）}

\subsection*{8-3}
实验测得反应的表观活化能在 $\cdot p ( \mathrm { O } _ { 2 } )$ 较⾼时为 72 kJ $\mathrm { m o l ^ { - 1 } }$ $p ( \mathrm { O } _ { 2 } )$ 较低时为 90kJ $\mathrm { m o l ^ { - 1 } }$ 。根据8-2中推导出的两种速率⽅程，指出两种情况下，分别是哪步的活化能决定了总反应的表观活化能？（提⽰：⾃由基相互结合的反应活化能通常近似为0）

\subsection*{8-4}
某种酶对上述反应的催化速率常数�随温度变化的关系如右表所⽰。其中，下标中的H或D分别表⽰8-1中亚油酸⼄酯上相应氢原⼦均是氕(H)或氘(D)原⼦。计算氕(H)原⼦情况下的表观活化能。

\textit{（原卷此处含表格/仪器试剂表，详见 PDF 或 MD 源文件。）}

\subsection*{8-5}
动⼒学同位素效应(KIE)指的是把反应物分⼦中的某⼀原⼦⽤其同位素原⼦取代后，反应速率常数变化的效应。⼀般⽤取代前后的速率常数之⽐衡量KIE的⼤⼩。计算8-4中各温度下的 $\mathrm { \ell } _ { \mathrm { t H } } / k _ { \mathrm { D } \circ }$ 并选择影响 $| k _ { \mathrm { H } } / k _ { \mathrm { D } }$ 数值的主要因素。（直接在答题纸上填A或B）A. 指前因⼦ B. 活化能

有机化学试题中的缩写

$\mathrm { T f _ { 2 } O } \colon$ 三氟甲磺酸酐 2-I-py：2-碘吡啶 4Bu：叔丁基

DABCO：1,4-⼆氮杂⼆环[2.2.2]⾟烷 'Pr：异丙基 \$Pr：正丙基

Ph：苯基 Cp：环戊⼆烯负离⼦ AcOH：⼄酸

THF：四氢呋喃 TsCl：对甲苯磺酰氯 DMAP：4-⼆甲氨基吡啶

LiHMDS：六甲基⼆硅氨基锂 AIBN：偶氮⼆异丁腈

EDCI：1-⼄基-(3-⼆甲基氨基丙基)碳⼆亚胺盐酸盐
